% Da migliorare per quanto riguarda l'umanità del testo

\section{Contesto}\label{sec:contesto}

La crescente disponibilità di dati relazionali da studiare negli ultimi anni ha portato ad un sempre maggiore interesse nello studio di strutture dati capaci di rappresentare interazioni complesse tra entità.
In letteratura, queste relazioni erano originariamente gestite mediante l'uso di \emph{grafi}, ovvero strutture matematiche utilizzate per modellare relazioni uno-a-uno tra oggetti.
Tuttavia, in molti scenari del mondo reale questa rappresentazione risulta troppo restrittiva, specialmente quando si considerano interazioni che coinvolgono più di due sole entità.
Diversi scenari della vita reale non possono essere rappresentati come grafi senza una significativa perdita di informazione.
Alcuni esempi includono le conversazioni tra più persone, che si trovano comunemente online nei forum, nelle chat o più in generale nei social media, così come le interazioni congiunte tra molteplici oggetti o entità, come quelle coinvolte in una transazione bancaria.
In questi contesti, la perdita di informazione è inevitabile se si limita la rappresentazione ad archi tra coppie di nodi.
Per superare queste limitazioni è stato introdotto il concetto di \emph{ipergrafo}.
Gli ipergrafi generalizzano i grafi standard consentendo a ciascun iperarco di connettere un numero arbitrario di nodi, rendendoli adatti a modellare interazioni di ordine superiore~\cite{antelmi2023}~\cite{bretto2013}.
Questo tipo di struttura dati è particolarmente utile quando l'obiettivo è rappresentare interazioni complesse senza perdere informazioni strutturali, che i grafi tradizionali tendono invece a scartare.
Inoltre, negli scenari pratici le interazioni tra entità non sono statiche, ma evolvono nel tempo.
Nuove interazioni vengono create, mentre altre tendono a scomparire, dando vita a un sistema complesso e dinamico che porta una struttura dati teoricamente statica a evolversi nel tempo.
Questo tipo di sfida rende necessaria la gestione di tale evoluzione temporale, al fine di poter rappresentare i dati sempre in modo accurato.
``Giusto per citare alcuni esempi, la modellazione di grafi temporali e l'analisi delle proprietà temporali possono trovare applicazione nella sociologia e nell'analisi delle reti sociali [\ldots]; nella sicurezza e nel calcolo distribuito [\ldots]; nella biologia''~\cite{lotito2020}.
L'analisi di grandi strutture dinamiche introduce inoltre importanti problemi computazionali~\cite{lotito2020}.
In particolare, è spesso necessario aggiornare continuamente le informazioni di interesse senza poter ricalcolare da capo l'intera soluzione a ogni passo temporale.
Di conseguenza, lo studio di tecniche efficienti per il mantenimento di proprietà e strutture combinatorie in ambienti dinamici e temporali svolge un ruolo centrale.